\documentclass[12pt,a4paper]{article}
\usepackage[utf8]{inputenc}
\usepackage[T1]{fontenc}
\usepackage[ngerman]{babel}
\usepackage{amsmath,amssymb,amsthm}
\usepackage{geometry}
\usepackage{hyperref}
\usepackage{enumitem}
\usepackage{mathrsfs}
\usepackage{yfonts}
\usepackage{fancyhdr}
\geometry{margin=2.5cm}

\theoremstyle{plain}
\newtheorem{satz}{Satz}[section]
\newtheorem{theorem}{Theorem}[section]
\newtheorem{lemma}{Hilfssatz}[section]
\newtheorem{hilfssatz}{Hilfssatz}[section]

\theoremstyle{definition}
\newtheorem{definition}{Definition}[section]

\theoremstyle{remark}
\newtheorem{bemerkung}{Bemerkung}[section]

\pagestyle{fancy}
\fancyhf{}
\fancyhead[C]{\small \nouppercase{\leftmark}}
\fancyfoot[C]{\thepage}
\renewcommand{\headrulewidth}{0.4pt}

\title{Lehrbuch der Algebra --- Erster Band, Teil 4}
\author{Heinrich Weber}
\date{}

\begin{document}
\maketitle

\section*{Fünfter Abschnitt}
\addcontentsline{toc}{section}{Fünfter Abschnitt}

\subsection*{\S. 32. Satz von Frobenius.}
\addcontentsline{toc}{subsection}{\S. 32. Satz von Frobenius.}

Frobenius hat einen Satz aufgestellt, der als ein Gegenstück zum vierten Sylow'schen betrachtet werden kann, der sich so aussprechen lässt\footnote{Frobenius, Sitzungsberichte der Berl. Akademie, 4. Mai 1893.}:

\begin{satz}
Jede Gruppe, deren Grad ein Product von lauter verschiedenen Primzahlen ist, ist metacyklisch.
\end{satz}

Wir setzen also voraus, dass der Grad $n$ einer Gruppe $P$ durch keine Quadratzahl theilbar sei. Ist $t$ der grösste Primtheiler von $n$, so giebt es nach dem Cauchy-Sylow'schen Satze (\S. 29) einen Theiler $T$ vom Grade $t$ von $P$, der, weil $t$ eine Primzahl ist, eine cyklische Gruppe ist, also aus den Elementen
\[
T = 1, a, a^{2}, \ldots a^{t-1}
\]
besteht. Wenn sich nun nachweisen liesse, dass unter den gemachten Voraussetzungen $T$ ein Normaltheiler von $P$ ist, so könnte man dieselbe Schlussweise auf die Gruppe $P/T$, deren Grad $n:t$ ist, anwenden, und würde, wenn $t_{1}$ die zweitgrösste in $n$ aufgehende Primzahl ist, auf einen Normaltheiler $N$ der Gruppe $P/T$ vom Grade $t_{1}$ schliessen. Daraus würde dann aber folgen (\S. 6, 2.), dass $P$ einen Normaltheiler $T_{1}$ vom Grade $tt_{1}$ hat, von dem $T$ selbst wieder Normaltheiler ist. Indem man in gleicher Weise die Gruppe $P/T_{1}$ betrachtet u.~s.~f., ergiebt sich eine Compositionsreihe von $P$:
\[
P,\ T_{k},\ T_{k-1},\ \ldots\ T_{1},\ T,\ 1,
\]
in der die Indexreihe aus den der Grösse nach aufsteigend geordneten Primfactoren von $n$ besteht, und $P$ ergiebt sich als metacyklisch. Alles kommt also darauf an, nachzuweisen, dass $T$ ein Normaltheiler von $P$ ist.

Dieser Satz ist leichter zu beweisen, wenn man ihn als speciellen Fall eines allgemeineren betrachtet, als wenn man ihn direct angreift. Dieser allgemeinere Satz lautet unter den über $P$ und $n$ gemachten Voraussetzungen:

\begin{enumerate}[label=\alph*)]
\item Ist $n = \mu\nu$ in zwei Factoren zerlegt und ist jeder Primtheiler von $\nu$ grösser als jeder Primtheiler von $\mu$, so giebt es in $P$ \emph{einen} und \emph{nur einen} Theiler vom Grade $\mu$; oder auch:
\item Es giebt in $P$ ein Element, dessen Grad durch keinen Primtheiler von $\nu$ theilbar ist, und ein Element, dessen Grad nicht durch $p$ theilbar ist, wenn $p$ ein Primtheiler von $\mu$ ist.
\end{enumerate}

Wir beweisen diesen Satz durch vollständige Induction und nehmen denselben als bewiesen an für alle Gruppen, deren Grad kleiner als $n$ ist.

Es sei $p$ ein Primtheiler von $\mu$, und $p^{\alpha}$ die höchste in $\mu$ aufgehende Potenz von $p$, und es sei $U$ die Gesammtheit aller Elemente von $P$, deren Grad durch keine höhere Potenz von $p$ theilbar ist als $p^{\alpha-1}$. Dann ist die Anzahl der Elemente in $U$, deren Grad durch $p^{\alpha-1}$ theilbar ist, nach dem Theorem I, \S. 31 durch $p^{\alpha-1}$ theilbar, und daher auch durch $p$ theilbar. Es giebt also in $U$ Elemente, deren Grad nicht durch $p^{\alpha-1}$ theilbar ist.

Ist nun $b$ ein solches Element und $\beta$ sein Grad, so enthält die cyklische Gruppe $B$, die aus den Potenzen von $b$ besteht, $p^{\alpha-1}\beta$ Elemente von $U$, nämlich die $p^{\alpha-1}$ Potenzen von $b$, deren Exponenten nicht durch $p$ theilbar sind. Die Anzahl der Elemente von $U$, deren Grad durch $p^{\alpha-1}$ theilbar ist, ist also ein Vielfaches von $p^{\alpha-1}\beta$, und daher durch $p^{\alpha}\beta$ theilbar. Da nun $\beta$ nicht durch $p$ theilbar ist, so ist die Anzahl dieser Elemente durch $p^{\alpha}$ theilbar.

Da die Anzahl aller Elemente von $U$, deren Grad durch $p^{\alpha-1}$ theilbar ist, gleich $p^{\alpha-1}\nu$ ist, so folgt, dass $p^{\alpha-1}\nu$ durch $p^{\alpha}\beta$ theilbar ist, und daher $\nu$ durch $p\beta$ theilbar ist. Da aber $\beta$ nicht durch $p$ theilbar ist, so muss $\nu$ durch $p$ theilbar sein, was der Voraussetzung widerspricht, dass jeder Primtheiler von $\nu$ grösser als jeder Primtheiler von $\mu$ ist.

Es sei nun $p$ ein Primtheiler von $\mu$ und $P$ eine Gruppe vom Grade $n$, deren Grad durch keine Quadratzahl theilbar ist. Dann ist die Anzahl der Elemente von $P$, deren Grad durch $p$ theilbar ist, gleich $jp$, worin $j$ eine ganze Zahl ist. Da nun die Anzahl der Elemente, deren Grad durch $p$ theilbar ist, gleich der Anzahl der Elemente ist, deren Grad nicht durch $p$ theilbar ist, so folgt, dass $jp = n - jp$, also $n = 2jp$ ist. Da aber $n$ durch keine Quadratzahl theilbar ist, so muss $j = 1$ sein, und daher $n = 2p$. Dies ist aber nur möglich, wenn $n = 2$ oder $n = 6$ ist.

Wenn wir nun noch beweisen können, dass es in $U$ genau $(p-1)\nu$ Elemente giebt, deren Grad durch $p$ theilbar ist, so ist der Satz bewiesen. Denn dann giebt es in $U$ genau $\nu$ Elemente, deren Grad nicht durch $p$ theilbar ist, und diese bilden eine Gruppe vom Grade $\nu$.

Die Anzahl der Elemente von $U$, deren Grad durch $p$ theilbar ist, ist nach dem Theorem I, \S. 31 gleich $jp$, worin $j$ eine ganze Zahl ist. Da nun die Anzahl der Elemente, deren Grad durch $p$ theilbar ist, gleich der Anzahl der Elemente ist, deren Grad nicht durch $p$ theilbar ist, so folgt, dass $jp = p^{\alpha}\nu - jp$, also $p^{\alpha}\nu = 2jp$ ist. Da aber $p^{\alpha-1}\nu$ nicht durch $p$ theilbar ist, so muss $j = p^{\alpha-1}\nu$ sein, und daher giebt es in $U$ genau $(p-1)p^{\alpha-1}\nu$ Elemente, deren Grad durch $p$ theilbar ist.

Ist, wie vorhin, $k$ der Grad von $\nu$, und setzen wir, wenn $h$ durch $p$ theilbar ist, $h = p^{\alpha}h'$, und wenn $h$ nicht durch $p$ theilbar ist, $h = h'$, so ist die Anzahl der Elemente von $U$, deren Grad durch $p$ theilbar ist, gleich
\[
\sum_{h} h = p^{\alpha} \sum_{h'} h' = p^{\alpha}\nu.
\]
Da nun die Anzahl der Elemente von $U$, deren Grad durch $p^{\alpha}$ theilbar ist, gleich $p^{\alpha}\nu$ ist, so folgt, dass die Anzahl der Elemente, deren Grad durch $p^{\alpha}$ aber nicht durch $p^{\alpha+1}$ theilbar ist, gleich $(p-1)p^{\alpha-1}\nu$ ist.

Damit ist also der Satz von Frobenius vollständig bewiesen.


\subsection*{\S. 33. Gruppen vom Grade $p^{\alpha}q$.}
\addcontentsline{toc}{subsection}{\S. 33. Gruppen vom Grade $p^{\alpha}q$.}

Wir wenden uns nun zu dem allgemeineren Problem, Gruppen vom Grade $p^{\alpha}q$ zu untersuchen, wobei $p$ und $q$ Primzahlen sind. Wir setzen voraus, dass $p < q$ ist und dass $\alpha$ eine beliebige positive ganze Zahl ist.

Es sei $P$ eine Gruppe vom Grade $p^{\alpha}q$ und $Q$ ein Theiler von $P$ vom Grade $p^{\beta}q$, wobei $0 \leqq \beta < \alpha$ vorausgesetzt ist, und es sei $\beta$ so gross als möglich angenommen. Dann ist wegen des zweiten Sylow'schen Satzes (\S. 30) die Anzahl der conjugirten Gruppen $Q$, $Q'$, $Q''$, $\ldots$ $Q^{(j-1)}$ in $P$ gleich $p^{\alpha-\beta}$, und $j$ ist nicht durch $p$ theilbar.

Nun ist aber, wenn $R$ der grösste gemeinsame Theiler von $Q$ und $Q'$ ist, der Grad von $R$ gleich $p^{\beta}q/p^{\alpha-\beta} = p^{2\beta-\alpha}q$, und da $R$ ein Theiler von $Q$ ist, so muss $2\beta - \alpha \geqq 0$ sein, also $\beta \geqq \alpha/2$.

Es lässt sich hiernach der Satz beweisen:

\begin{satz}
$\gamma$) Jedes Element von $P$, dessen Grad durch $p$ theilbar ist, kommt in einer und nur in einer der Gruppen $R$, $R'$, $R''$, $\ldots$ $R^{(j-1)}$ vor.
\end{satz}

Die Gruppen $Q$, $Q'$, $Q''$, $\ldots$ sind, da ihr Grad eine Primzahl ist, cyklisch, und je zwei enthalten, weil sie von einander verschieden sind, ausser dem Einheitselemente kein gemeinsames Element. Es sei
\[
Q = 1, a, a^{2}, \ldots a^{p-1}.
\]

Die Gruppe $R$, die aus allen der Bedingung
\[
c^{-1} Q c = Q
\]
genügenden Elementen besteht, enthält ausser den Potenzen von $a$ kein Element vom Grade $p$; und Entsprechendes gilt für die anderen Gruppen $Q'$, $R'$, $Q''$, $R''$, $\ldots$

Ist nun $b$ ein Element von $P$, dessen Grad durch $p$ theilbar ist und gleich $ps$ sein mag, so dass $s$ nicht durch $p$ theilbar ist, so ist $b^{s}$ ein Element, dessen Grad $p$ ist, und das also in einer und nur in einer der cyklischen Gruppen $Q$, $Q'$, $Q''$, $\ldots$ vorkommt. Ist aber $b^{s} = a$ ein Element von $Q$, so ist
\[
b^{-1} a b = a,
\]
d.~h. $b$ kommt in der Gruppe $R$ vor, und kann in keiner der anderen Gruppen, z.~B. in $R'$, vorkommen, weil sonst $b^{s}$ auch in $Q'$ vorkäme, was nicht möglich ist, da $Q$ und $Q'$ nur das Einheitselement gemein haben. Damit ist $\gamma$) bewiesen.

\begin{satz}
$\delta$) In $U$ giebt es genau $jr_{2}(p-1)$ Elemente, deren Grad durch $p$ theilbar ist.
\end{satz}

Unter $U$ haben wir die Gesammtheit der Elemente von $P$ verstanden, deren Grad ein Theiler von $pr$ ist. Nun giebt es nach $\beta$) in $R$ genau $r_{2}(p-1)$ Elemente aus $U$, deren Grad durch $p$ theilbar ist, und da jede der $j$ Gruppen $R$, $R'$, $R''$, $\ldots$ an Stelle von $R$ treten kann, so folgt $\delta$) aus $\gamma$). Um also noch zu zeigen, worauf nach dem Obigen alles ankommt, dass in $U$ genau $j r_{2}(p-1)$ Elemente vorkommen, deren Grad durch $p$ theilbar ist, ist also nur noch die Formel
\[
jp = jr_{2}(p-1)
\]
zu beweisen, die aber unmittelbar aus der Definition von $r_{2}$ folgt.

Damit ist also das Theorem II vollständig bewiesen.

Nimmt man in diesem Satze $\beta = \alpha - 1$ an, so fällt $R$ mit $P$ zusammen, und es ergiebt sich, dass $Q$ ein Normaltheiler von $P$ ist. Wendet man denselben Satz auf die Gruppe $Q$ an u.~s.~f., so erhält man:

\begin{satz}
IV. Jede Gruppe, deren Grad eine Primzahlpotenz ist, ist metacyklisch (\S. 8).
\end{satz}


\subsection*{\S. 34. Einfache Gruppen.}
\addcontentsline{toc}{subsection}{\S. 34. Einfache Gruppen.}

Die Untersuchung der einfachen Gruppen ist eines der wichtigsten und schwierigsten Probleme der Gruppentheorie. Wir haben bereits gesehen, dass die Gruppen, deren Grad eine Primzahl oder eine Primzahlpotenz ist, stets metacyklisch sind und daher keine einfachen Gruppen sein können, ausser wenn sie selbst Primzahlengrades sind.

Wir wollen nun die einfachen Gruppen von niederem Grade untersuchen und feststellen, welche Grade eine einfache Gruppe überhaupt haben kann. Wir haben bereits gesehen, dass die einzigen einfachen Gruppen von Primzahlengrad die cyklischen Gruppen von Primzahlengrad sind, und dass es keine einfachen Gruppen von Primzahlpotenzgrad giebt, ausser den cyklischen.

Es bleibt nur noch die Untersuchung der Zahlen 60 und 90 übrig. Dass eine einfache Gruppe vom 60\textsuperscript{ten} Grade existirt, wissen wir schon; nämlich die alternirende Gruppe der Permutationen von fünf Buchstaben (Bd.~I, \S.~177). Es ist aber noch fraglich, ob dies die einzige ist.

Eine einfache Gruppe $P$ vom Grade 60 enthält nach \S.~29, I. als Theiler eine Gruppe 5\textsuperscript{ten} Grades und eine Gruppe 3\textsuperscript{ten} Grades. Nimmt man in dem zweiten Sylow'schen Satze (\S.~30) für $Q$ die Gruppe 5\textsuperscript{ten} Grades, so muss $R$ vom 10\textsuperscript{ten} Grade sein, weil der Index von $R$ congruent mit 1 nach dem Modul 5, also nur $= 6$ sein kann. Also kann die Gruppe $P$ vom 60\textsuperscript{ten} Grade nach \S.~28, 2. als transitive Permutationsgruppe von sechs Ziffern dargestellt werden; sie kann aber keine cyklische Permutation von nur drei Ziffern enthalten, weil sie als einfache Gruppe (nach Bd.~I, \S.~158, 2.) primitiv sein muss, und daher, wenn sie einen dreigliedrigen Cyklus enthielte, nach Bd.~I, \S.~153, 10. die ganze alternirende Gruppe 360\textsuperscript{ten} Grades enthalten müsste. Wir können also eine der Permutationen 3\textsuperscript{ten} Grades in der Form annehmen:
\[
a = (1,2,3)(4,5,6).
\]

Hierzu wollen wir eine Permutation 5\textsuperscript{ten} Grades, $b$, fügen, die nur eine cyklische sein kann. Lassen wir darin die Ziffer 6 fehlen, so können wir statt $b$ eine solche Potenz von $b$ nehmen (die ja auch in $P$ vorkommen muss), dass etwa 1, 2 die beiden ersten Ziffern werden, und dann bleiben noch sechs Möglichkeiten für $b$ übrig:
\begin{align*}
&(1,2,3,4,5),\ (1,2,4,3,5),\ (1,2,4,5,3),\\
&(1,2,3,5,4),\ (1,2,5,3,4),\ (1,2,5,4,3);
\end{align*}
von diesen sechs Annahmen sind aber wieder die drei in einer Reihe stehenden nicht wesentlich (d.~h. nur durch die Bezeichnung) verschieden. Denn ersetzt man z.~B. bei der zweiten Annahme $b = (1,2,4,3,5)$ das Element $a$ durch $a^{2}$ und $b$ durch $b^{3}$, so erhält man
\[
(1,3,2)(4,6,5);\ (1,3,2,5,4),
\]
was durch Vertauschung der Ziffern 2 mit 3 und 4 mit 5 in die erste Annahme übergeht. Die dritte Annahme $b = (1,2,4,5,3)$ geht, wenn man $b$ durch $b^{4}$ und $a$ durch $a^{2}$ ersetzt und dann 1 mit 2 und 4 mit 5 vertauscht, in die erste über, und ebenso lassen sich die drei anderen Annahmen über $b$ auf einander zurückführen.

Dass eine einfache Permutationsgruppe 60\textsuperscript{ten} Grades bei fünf Ziffern wirklich existirt, nämlich die alternirende Gruppe, haben wir schon früher (Bd.~I, \S.~177) nachgewiesen; und es ergiebt sich also, dass diese bis auf Isomorphismus die einzige einfache Gruppe vom Grade 60 ist.


\subsection*{\S. 35. Gruppen vom Grade $pq$.}
\addcontentsline{toc}{subsection}{\S. 35. Gruppen vom Grade $pq$.}

Wir gehen nun zu dem Falle über, dass der Grad $pq$ das Product zweier Primzahlen ist, die auch einander gleich sein können. Ist $p = q$, so ist der Grad $p^{2}$, und die Gruppe ist nach dem vorhergehenden Paragraph metacyklisch.

Wir können also annehmen, dass $p$ und $q$ verschiedene Primzahlen sind. Ist $p < q$, so enthält die Gruppe $P$ vom Grade $pq$ nach dem Cauchy-Sylow'schen Satze einen Theiler vom Grade $q$, und die Anzahl dieser Theiler ist congruent 1 nach dem Modul $q$, also gleich 1, wenn $p$ nicht congruent 1 nach dem Modul $q$ ist. In diesem Falle ist also der Theiler vom Grade $q$ ein Normaltheiler, und die Gruppe ist metacyklisch.

Ist aber $p$ congruent 1 nach dem Modul $q$, so kann die Gruppe nicht metacyklisch sein. Es giebt dann nämlich mehrere verschiedene Theiler vom Grade $q$, und die Gruppe enthält Elemente vom Grade $p$, die nicht in einem Normaltheiler enthalten sind.

Die Gruppe $P$ enthält in diesem Falle einen Theiler $Q$ vom Grade $q$ und einen Theiler $R$ vom Grade $p$, und es ist $Q$ ein Normaltheiler von $P$, wenn $q$ nicht congruent 1 nach dem Modul $p$ ist. Ist aber $q$ congruent 1 nach dem Modul $p$, so ist $Q$ kein Normaltheiler, und die Gruppe ist nicht metacyklisch.


\subsection*{\S. 36. Grenzen des Index eines Theilers der symmetrischen Permutationsgruppe.}
\addcontentsline{toc}{subsection}{\S. 36. Grenzen des Index eines Theilers der symmetrischen Permutationsgruppe.}

Wir beschliessen diese Betrachtungen mit dem Beweise eines Satzes über Permutationsgruppen, der durch die Schwierigkeit, die sein Beweis anfangs bot, eine gewisse Berühmtheit erlangt hat, und der für die Beurtheilung algebraischer Fragen von Wichtigkeit ist\footnote{Jordan, Traité des substitutions.}.

Es handelt sich dabei um die symmetrische Permutationsgruppe $P$ von $n$ Ziffern und um ihre Theiler von möglichst kleinem Index. Wir wissen, dass die Gruppe $P$ immer einen Theiler vom Index 2 hat, nämlich die alternirende Gruppe. Ausserdem ist noch ein Theiler vom Index $n$ bekannt, der alle Permutationen von $P$ umfasst, die eine Ziffer ungeändert lassen, der also intransitiv ist.

Zunächst gilt der folgende Satz:

\begin{satz}
I. Der Index eines imprimitiven Theilers von $P$ ist immer grösser als $n$, und der Index eines intransitiven Theilers ist gleich oder grösser als $n$, und nur dann gleich $n$, wenn der Theiler eine Ziffer in Ruhe lässt, und die übrigen $n-1$ Ziffern auf alle mögliche Arten permutirt.
\end{satz}

Nehmen wir an, es sei $Q$ ein imprimitiver Theiler von $P$ vom Index $j$, und es bestehen $r$ Systeme der Imprimitivität von je $s$ Ziffern, so dass $n = rs$ ist. Eine Zahl, die der Grad der Gruppe $Q$ sicher nicht übersteigen kann, erhalten wir, wenn wir alle Permutationen in jedem einzelnen der $r$ Systeme und dann noch sämmtliche Permutationen der Systeme abzählen. Der Grad von $Q$ ist also kleiner oder gleich
\[
(s!)^{r} \cdot r! = \frac{(s!)^{r} \cdot r!}{rs} \cdot n.
\]
Dies ist aber kleiner als $n/j$, also ist $j > n$, wenn nicht $s = 1$ oder $r = 1$ ist, was aber den Voraussetzungen widerspricht.

Für einen intransitiven Theiler ergiebt sich der Satz unmittelbar aus der Definition.

Wir nehmen nun an, es sei $Q$ ein primitiver Theiler von $P$ vom Index $j$, und es sei $b$ eine Permutation von $Q$, die die Ziffer $a$ in $b$ verwandelt. Dann ist der Complex $Pb$ mit $P$ völlig identisch. Es können sich also die beiden Reihen
\begin{align*}
A &= a_{0},\ a_{1},\ a_{2},\ \ldots,\ a_{n-1}\\
Ab &= a_{0}b,\ a_{1}b,\ a_{2}b,\ \ldots,\ a_{n-1}b
\end{align*}
nur durch die Anordnung von einander unterscheiden. Der Uebergang von $A$ zu $Ab$ ist also eine Permutation von $n$ Ziffern $0, 1, 2, \ldots n-1$, und jedem Elemente $b$ von $P$ entspricht eine solche Permutation, die wir mit $\pi_{b}$ bezeichnen wollen. Zwei verschiedenen Elementen $b$, $c$ entsprechen immer zwei verschiedene Permutationen $\pi_{b}$, $\pi_{c}$, und dem Einheitselemente entspricht die identische Permutation. Setzen wir
\[
\pi_{b} = (A, Ab),\quad \pi_{c} = (A, Ac),
\]
so können wir $\pi_{c}$ auch mit $(Ab, Abc)$ bezeichnen, und daraus ergiebt sich
\begin{equation}
\pi_{b}\, \pi_{c} = \pi_{bc},
\end{equation}
d.~h. die Permutationen $\pi$ bilden eine mit $P$ isomorphe Permutationsgruppe. Diese Permutationsgruppe ist transitiv, da z.~B. das Element $a_{0}$ in jedes beliebige andere Element von $P$ übergehen kann. Also haben wir den Satz:

\begin{satz}
1. Jede Gruppe vom Grade $n$ ist isomorph mit einer transitiven Permutationsgruppe von $n$ Ziffern.
\end{satz}

Es giebt natürlich noch andere mit einer gegebenen Gruppe isomorphe Permutationsgruppen, und es wäre von besonderem Interesse, eine solche Gruppe mit möglichst geringer Ziffernzahl zu bilden. Diese Aufgabe kann bis jetzt nicht allgemein gelöst werden. Wir müssen uns hier mit wenigen Sätzen begnügen.

Wir nehmen an, es sei $R$ irgend ein Divisor von $P$ vom Index $j$, und bezeichnen die Elemente von $R$ mit $c$, setzen ferner, indem wir $P$ in ein System von Nebengruppen zerlegen,
\begin{equation}
P = R + Rb_{1} + Rb_{2} + \cdots + Rb_{j-1}.
\end{equation}
Ist dann $a$ irgend ein Element von $P$, so werden die beiden Systeme
\begin{align}
B &= R,\ Rb_{1},\ Rb_{2},\ \ldots,\ Rb_{j-1}\\
Ba &= Ra,\ Rb_{1}a,\ Rb_{2}a,\ \ldots,\ Rb_{j-1}a,\nonumber
\end{align}
von der Reihenfolge abgesehen, übereinstimmen, und es entspricht also jedem Element $a$ eine bestimmte Permutation $\pi_{a}$ der $j$ Ziffern $0, 1, \ldots j-1$, die wir mit
\begin{equation}
\pi_{a} = (B, Ba)
\end{equation}
bezeichnen.

Ist nun $a$ in diesen $j-1$ Ziffern intransitiv, so müsste nach (1) der Grad $\mu_{a}$ von $\pi_{a}$ kleiner als $j$ sein. Wenn aber durch $Q_{0}: a$ noch eine vierte Ziffer 3 ungeändert bliebe, so würde $a$ in den Ziffern 0, 1, 2, 3 intransitiv sein, und der Grad von $\pi_{a}$ wäre höchstens $j-1$. Wenn wir nun annehmen, dass durch $Q_{0}: a, b$ noch eine fünfte Ziffer 4 ungeändert bliebe, so würde der Grad von $\pi_{a}$ höchstens $j-2$ sein, u.~s.~f.

Der Grad $\mu_{a}$ ergiebt sich genau wie der von $Q$ und $Q_{0}$: wenn wir annehmen, dass $a$ in den Ziffern 0, 1, 2 intransitiv ist, so ist $\mu_{a} = j - 1$; wenn $a$ in 0, 1, 2, 3 intransitiv ist, so ist $\mu_{a} = j - 2$, u.~s.~f. Allgemein, wenn $a$ in $k+1$ Ziffern intransitiv ist, so ist $\mu_{a} = j - k$.

Wir nehmen, was freisteht, an, dass $b > a$ sei, so ist der Ausdruck auf der rechten Seite dieser Ungleichung grösser als
\[
\frac{1}{2}\,\frac{n-1}{n-a}\,\frac{n-a}{n-b} = \frac{1}{2}\,\frac{n-1}{n-b},
\]
und da $b \leqq n/2$ ist, so ist dies grösser als 1. Also ist in der That der Grad von $Q_{0}$ grösser als $n$, und der Index von $Q$ ist grösser als $n$.

Damit ist also der Satz bewiesen:

\begin{satz}
II. Der Index eines primitiven Theilers der symmetrischen Gruppe von $n$ Ziffern ist, wenn er von der Gruppe selbst und der alternirenden Gruppe verschieden ist, immer grösser als $n$, mit Ausnahme der Fälle $n = 6$ und $n = 12$.
\end{satz}

Für $n = 6$ existirt eine Gruppe vom Grade 120, die einen Theiler vom Index 6 hat, nämlich die symmetrische Gruppe selbst; und für $n = 12$ existirt eine Gruppe vom Grade 95040 (die Mathieu'sche Gruppe), die primitive Theiler von viel niederem Index als $n$ besitzt.


\newpage
\section*{Zweites Buch. Lineare Gruppen.}
\addcontentsline{toc}{section}{Zweites Buch. Lineare Gruppen.}

\section*{Sechster Abschnitt. Gruppen linearer Substitutionen.}
\addcontentsline{toc}{section}{Sechster Abschnitt. Gruppen linearer Substitutionen.}

\subsection*{\S. 37. Lineare Substitutionen.}
\addcontentsline{toc}{subsection}{\S. 37. Lineare Substitutionen.}

Wir wenden uns nun zu einer wichtigen und umfangreichen Klasse von Gruppen, den Gruppen linearer Substitutionen. Unter einer \emph{linearen Substitution} verstehen wir eine Substitution der Form
\begin{equation}
x_{i}' = a_{i1}x_{1} + a_{i2}x_{2} + \cdots + a_{in}x_{n}\quad (i = 1, 2, \ldots, n),
\end{equation}
worin die $a_{ik}$ Constanten sind, die einer bestimmten Determinante
\[
|A| = |a_{ik}|\quad (i,k = 1, 2, \ldots, n)
\]
bedürfen. Die $n^{2}$ Grössen $a_{ik}$ nennen wir die \emph{Coefficienten} oder \emph{Elemente} der Substitution. Die Zahl $n$ heisst die \emph{Dimension} oder der \emph{Grad} der Substitution.

Zur bequemeren Bezeichnung fassen wir die Coefficienten $a_{ik}$ in eine quadratische Matrix zusammen, die wir mit $(A)$ oder einfach mit $A$ bezeichnen, und schreiben die Substitution in der Form
\begin{equation}
(x') = A(x),
\end{equation}
worin $(x)$ das System der $n$ Variablen $x_{1}, x_{2}, \ldots, x_{n}$ und $(x')$ das System der $n$ neuen Variablen $x_{1}', x_{2}', \ldots, x_{n}'$ bedeutet.

\emph{Die identische Substitution} ist die Substitution
\[
x_{i}' = x_{i}\quad (i = 1, 2, \ldots, n),
\]
deren Coefficienten die Einheitsmatrix
\[
E = \begin{pmatrix}
1 & 0 & \cdots & 0\\
0 & 1 & \cdots & 0\\
\vdots & \vdots & \ddots & \vdots\\
0 & 0 & \cdots & 1
\end{pmatrix}
\]
bilden.

Sind zwei lineare Substitutionen $A$ und $B$ gegeben:
\begin{equation}
(y) = A(x),\quad (z) = B(y),
\end{equation}
so können wir die Variablen $z$ unmittelbar durch die $x$ ausgedrückt mittelst einer neuen linearen Substitution $E$, die wir so bezeichnen können:
\begin{equation}
(z) = E(x) = BA(x),
\end{equation}
und die Substitution $E = BA$ heisst aus $B$ und $A$ \emph{zusammengesetzt} oder \emph{componirt}. Es ist dabei aber zwischen $AB$ und $BA$ zu unterscheiden. Zwei Substitutionen $A$, $B$ von der besonderen Eigenschaft, dass $AB = BA$ ist, heissen mit einander \emph{vertauschbar} oder \emph{commutativ}. Die Substitutionscoefficienten von $E$ ergeben sich durch Einsetzen der Ausdrücke (1) in (9):
\begin{equation}
e_{ik} = \sum_{h=1}^{n} b_{ih}a_{hk}.
\end{equation}

Diese Formeln sind ganz dieselben, die wir im \S.~27 des ersten Bandes benutzt haben\footnote{Vgl. Bd.~I, \S.~27.}, und wir können demnach das in (11) ausgedrückte Gesetz der Composition der Substitutionen so ausdrücken:

\begin{satz}
1. Um die aus zwei Substitutionen $B$, $A$ zusammengesetzte Substitution $BA$ zu bilden, verfährt man ganz so, als ob die beiden Determinanten $|B|$, $|A|$ nach der Multiplicationsregel mit einander multiplicirt werden sollten. Es sind dabei, um die Elemente einer Zeile zu bilden, die Elemente einer Zeile der ersten Componente mit den entsprechenden Elementen der Colonnen der zweiten Componente zu multipliciren und dann zu addiren. Man drückt dies auch kurz so aus, dass in der ersten Componente nach Zeilen, in der zweiten nach Colonnen summirt wird.
\end{satz}

Aus dieser Regel ergiebt sich die Folgerung:

\begin{satz}
2. Die Determinante einer zusammengesetzten Substitution ist gleich dem Producte aus den Determinanten der Componenten.
\end{satz}

Denn es ist
\[
|E| = |BA| = |B| \cdot |A|.
\]

\subsection*{\S. 37. Inverse Substitutionen.}
\addcontentsline{toc}{subsection}{\S. 37. Inverse Substitutionen.}

Ebenso leicht erhalten wir:

\begin{satz}
3. Durch Zusammensetzung mit der identischen Substitution ändert sich eine Substitution nicht.
\end{satz}

Denn es ist $EA = AE = A$.

\begin{satz}
4. Zu jeder Substitution $A$, deren Determinante von Null verschieden ist, giebt es eine und nur eine inverse Substitution $A^{-1}$, so dass
\[
AA^{-1} = A^{-1}A = E.
\]
\end{satz}

Die Coefficienten der inversen Substitution $A^{-1}$ sind die mit den richtigen Vorzeichen genommenen Unterdeterminanten von $|A|$, dividirt durch $|A|$ selbst. Ist also
\[
A = (a_{ik}),\quad A^{-1} = (\alpha_{ik}),
\]
so ist
\[
\alpha_{ik} = \frac{A_{ki}}{|A|},
\]
worin $A_{ki}$ die zu $a_{ki}$ gehörige Unterdeterminante bedeutet.

\begin{satz}
5. Die Gesammtheit aller linearen Substitutionen von bestimmter Dimension $n$ bildet eine (unendliche) Gruppe (\S.~1).
\end{satz}

Wenn wir aus dieser Gesammtheit irgend eine Menge herausgreifen, die die Eigenschaft hat, dass das Product von je zwei derselben angehörigen Substitutionen wieder der Menge angehört, und dass mit jeder Substitution auch ihre inverse in der Menge enthalten ist, so bilden diese Substitutionen eine \emph{endliche} oder \emph{unendliche Gruppe linearer Substitutionen}.

Die \emph{Ordnung} einer endlichen Gruppe linearer Substitutionen ist die Anzahl der in ihr enthaltenen Substitutionen.

\subsection*{\S. 37. Contragradiente Substitutionen.}
\addcontentsline{toc}{subsection}{\S. 37. Contragradiente Substitutionen.}

Neben der Substitution $A$ betrachten wir noch die \emph{transponirte} oder \emph{contragrediente Substitution} $A'$, die aus $A$ durch Vertauschung der Zeilen mit den Colonnen entsteht. Die Coefficienten von $A'$ sind also
\[
a'_{ik} = a_{ki}.
\]

Die contragrediente Substitution hat die Eigenschaft, dass für zwei Substitutionen $A$, $B$
\[
(AB)' = B'A'.
\]
Diese Bemerkung giebt die Vorschrift, nach der die transponirten Substitutionen zusammengesetzt werden, die sich in dem entgegengesetzten Sinne wie die ursprünglichen verhalten.


\subsection*{\S. 38. Substitution der Verhältnisse.}
\addcontentsline{toc}{subsection}{\S. 38. Substitution der Verhältnisse.}

Im Gebiete der binären Substitutionen ist die Substitution $A$:
\begin{equation}
\begin{pmatrix} x_{1}' \\ x_{2}' \end{pmatrix} = \begin{pmatrix} a_{11} & a_{12} \\ a_{21} & a_{22} \end{pmatrix} \begin{pmatrix} x_{1} \\ x_{2} \end{pmatrix}
\end{equation}
durch die Verhältnisse $z = x_{1}/x_{2}$ und $z' = x_{1}'/x_{2}'$ ausgedrückt:
\[
z' = \frac{a_{11}z + a_{12}}{a_{21}z + a_{22}}.
\]

Diese Formel bestimmt eine \emph{gebrochene lineare Substitution} der einen Variablen $z$, die als \emph{Substitution der Verhältnisse} bezeichnet wird.

Bezeichnen wir nämlich mit $z_{1}$, $z_{2}$, $\ldots$, $z_{n}$ ein System von $n$ Veränderlichen, und mit $0, 0', \ldots$ irgend eine Anordnung der $n+1$ Verhältnisse $z_{1}:z_{2}:\cdots:z_{n}:1$, so erhalten wir eine Substitution der Verhältnisse, wenn wir die lineare Substitution $A$ auf die homogenen Coordinaten anwenden.


\subsection*{\S. 40. Invarianten endlicher Gruppen.}
\addcontentsline{toc}{subsection}{\S. 40. Invarianten endlicher Gruppen.}

Wir betrachten eine endliche Gruppe $S$ linearer Substitutionen von $n$ Variablen $x_{1}, x_{2}, \ldots, x_{n}$. Eine \emph{Form} $F(x_{1}, x_{2}, \ldots, x_{n})$ heisst eine \emph{Invariante} oder \emph{absolute Invariante} der Gruppe $S$, wenn sie durch alle Substitutionen der Gruppe ungeändert bleibt, d.~h. wenn für jede Substitution $A$ der Gruppe
\[
F(Ax) = F(x).
\]

Eine Form $F(x)$ heisst eine \emph{relative Invariante} der Gruppe $S$, wenn sie durch jede Substitution $A$ der Gruppe bis auf einen Factor multiplicirt wird:
\[
F(Ax) = \chi(A) \cdot F(x),
\]
worin $\chi(A)$ eine von $A$ abhängige Grösse ist, die dem Multiplicationsgesetz
\[
\chi(AB) = \chi(A)\chi(B)
\]
genügt. Die Function $\chi(A)$ heisst der \emph{Charakter} der relativen Invariante.

Da $S$ eine endliche Gruppe ist, so muss $\chi(A)$ eine Einheitswurzel sein, und es soll der \emph{Index} der Invariante $F(x)$ heissen, die kleinste positive Zahl $\varrho$, für die $\chi(A)^{\varrho} = 1$ ist für alle $A$ in $S$. Ist $\varepsilon$ eine primitive $\varrho$\textsuperscript{te} Einheitswurzel, so können wir
\[
\chi(A) = \varepsilon^{\alpha(A)}
\]
setzen, worin $\alpha(A)$ eine ganze Zahl ist.

\begin{satz}
Ist der Charakter $\chi(A)$ nicht identisch gleich 1, so bilden die Substitutionen $A$, für die $\chi(A) = 1$ ist, eine Gruppe $T$, die ein Normaltheiler von $S$ ist, und der Index von $T$ in $S$ ist gleich dem Index $\varrho$ der Invarianten.
\end{satz}

Dies ergiebt sich einfach daraus, dass eine Transposition, z.~B. $(1,2)$, der Substitution
\[
F(x_{1}, x_{2}, \ldots, x_{n}) = x_{1} - x_{2}
\]
entspricht.

Hat die Gruppe $S$ vom Grade $\mu$ eine Invariante vom Index $\varrho$, so wird $T$ die Einheitsgruppe und $S$ ist eine cyklische Gruppe, deren Ordnung gleich dem Index der Invariante ist.


\subsection*{\S. 40. Absolute und relative Invarianten.}
\addcontentsline{toc}{subsection}{\S. 40. Absolute und relative Invarianten.}

Wir nehmen eine beliebige Form $F(x)$ der Veränderlichen $x$ zu der Gesammtheit der Functionen:
\[
F(Ax),\quad F(Bx),\quad F(Cx),\ \ldots
\]
indem wir für $A$, $B$, $C$, $\ldots$ alle Substitutionen der Gruppe $S$ setzen. Diese Functionen sind mit $F$ von gleichem Grade, und ihre Anzahl ist höchstens gleich der Ordnung $\mu$ von $S$. Bilden wir die symmetrischen Functionen dieser $\mu$ Functionen, so erhalten wir $\mu$ absolute Invarianten der Gruppe $S$, die mit $J_{1}$, $J_{2}$, $\ldots$, $J_{\mu}$ bezeichnet werden mögen.

Wir bezeichnen mit $\Omega$ den Rationalitätsbereich, der aus den absoluten Invarianten der Gruppe $S$ und allen rationalen Zahlen besteht. Dann ist jede absolute Invariante von $S$ eine rationale Function der $J_{1}$, $J_{2}$, $\ldots$, $J_{\mu}$ mit Coefficienten aus $\Omega$.


\subsection*{\S. 41. Satz von Hilbert.}
\addcontentsline{toc}{subsection}{\S. 41. Satz von Hilbert.}

Wir beweisen nun einen fundamentalen Satz über die Invarianten endlicher Gruppen, der auf Hilbert zurückgeht.

\begin{satz}
Das vollständige Invariantensystem einer endlichen Gruppe linearer Substitutionen ist stets \emph{endlich}, d.~h. es giebt immer eine endliche Anzahl von Invarianten $J_{1}$, $J_{2}$, $\ldots$, $J_{k}$, so dass jede andere Invariante der Gruppe sich als ganze rationale Function dieser $k$ Invarianten darstellen lässt.
\end{satz}

Dieser Satz ist ein specieller Fall eines allgemeineren Satzes von Hilbert über die Endlichkeit der Basis eines Systems von Formen.

Wir bezeichnen mit $\mathfrak{I}$ das System aller nicht constanten Invarianten der Gruppe $S$ und beweisen zunächst:

\begin{satz}
I. Es giebt in $\mathfrak{I}$ eine endliche Anzahl von Invarianten $J_{1}$, $J_{2}$, $\ldots$, $J_{m}$, so dass sich jede Invariante aus $\mathfrak{I}$ in der Form darstellen lässt:
\[
F = A_{1}J_{1} + A_{2}J_{2} + \cdots + A_{m}J_{m},
\]
worin $A_{1}$, $A_{2}$, $\ldots$, $A_{m}$ ganze Invarianten (eventuell Constanten) sind.
\end{satz}

Wir ordnen die Invarianten nach ihrem Grade und wählen aus den Invarianten vom niedrigsten Grade $r$ eine endliche Basis $F_{1}$, $F_{2}$, $\ldots$, $F_{\mu}$. Dann bilden wir die $\mu$ Producte $A_{1}F_{1}$, $A_{2}F_{2}$, $\ldots$, $A_{\mu}F_{\mu}$ alle von gleichem Grade, dem Grade von $F$. Natürlich aber wird im Allgemeinen nicht jede Invariante vom Grade $r$ sich als lineare Combination dieser Producte darstellen lassen.

Es sei $F_{1}$, $F_{2}$, $\ldots$, $F_{\nu}$ eine nach II bestimmte Basis des Systemes $\mathfrak{I}$, und $F$ irgend eine andere nicht constante Invariante. Wir ordnen $F$ nach Potenzen von $y$:
\[
F = \Phi_{0} + \Phi_{1}y + \Phi_{2}y^{2} + \cdots + \Phi_{s}y^{s},
\]
worin $\Phi_{0}$, $\Phi_{1}$, $\ldots$, $\Phi_{s}$ Formen der Variablen $z_{1}$, $z_{2}$, $\ldots$, $z_{n}$ sind. Aus (2), (3), (4) ergiebt sich aber, wenn wir noch
\[
P = b_{0} - aA_{1} - GA_{2} - GA_{3}
\]
setzen, eine Gleichung, die für alle Werthe der Variablen gültig ist.


\subsection*{\S. 42. Endlichkeit des Invariantensystemes einer endlichen linearen Substitutionsgruppe.}
\addcontentsline{toc}{subsection}{\S. 42. Endlichkeit des Invariantensystemes einer endlichen linearen Substitutionsgruppe.}

Der im vorhergehenden Paragraph allgemein bewiesene Satz von Hilbert lässt sich für den Fall einer endlichen Gruppe linearer Substitutionen noch wesentlich verschärfen. Wir zeigen nämlich, dass für eine endliche Gruppe das Invariantensystem nicht nur endlich ist, sondern dass sich sogar die Anzahl der zur Bildung der Basis nötigen Invarianten nach oben begrenzen lässt.

Es sei $F_{1}$, $F_{2}$, $\ldots$, $F_{\nu}$ eine nach II bestimmte Basis des Systemes $\mathfrak{I}$, und $F$ irgend eine andere nicht constante Invariante vom Grade $r$. Wir ordnen $F$ nach Potenzen von $y$ und führen die Division mit $F_{1}$ aus, wobei sich
\begin{equation}
F = a\Phi + \mathfrak{R}
\end{equation}
ergieben mag, so dass $a$ der Quotient und $\mathfrak{R}$ der Rest der Division ist. $a$ und $\mathfrak{R}$ sind ganze Functionen der Variablen $x$, $y$, weil der Coefficient der höchsten Potenz von $y$ im Divisor $F_{1}$ constant ist. $\mathfrak{R}$ übersteigt in Bezug auf $y$ nicht den Grad $r-1$.

Durchläuft nun $F$ das System $\mathfrak{I}$, so bildet die Gesammtheit der durch (11) definirten Functionen $\mathfrak{R}$ ein System $\mathfrak{I}_{r-1}$, von dem wir die Darstellbarkeit durch eine Basis als schon erwiesen annehmen. Wir können also setzen:
\begin{equation}
\mathfrak{R} = A_{2}\Phi_{2} + \cdots + A_{\nu}\Phi_{\nu},
\end{equation}
so dass $\Phi_{2}$, $\ldots$, $\Phi_{\nu}$ dem Systeme $\mathfrak{I}_{r-1}$ angehören, d.~h. so, dass sich in dem Systeme $\mathfrak{I}$ die Formen
\begin{equation}
F_{2} = a + \Phi_{2},\ \ldots,\ F_{\nu} = a + \Phi_{\nu}
\end{equation}
bestimmen lassen. Setzen wir also
\begin{equation}
A_{1} = a - aA_{2} - \cdots - aA_{\nu},
\end{equation}
so folgt aus (11), (12) und (13):
\begin{equation}
F = A_{1}F_{1} + A_{2}F_{2} + \cdots + A_{\nu}F_{\nu},
\end{equation}
wodurch das Theorem I. allgemein bewiesen ist.

Ferner das System $\mathfrak{I}$ Formen 0\textsuperscript{ten} Grades, d.~h. von Null verschiedene Constanten enthält, so ist, da wir für $F_{1}$ eine solche Form wählen können, das Theorem I. auch in diesem Falle gültig.


\subsection*{\S. 43. Das Formenproblem.}
\addcontentsline{toc}{subsection}{\S. 43. Das Formenproblem.}

Wir wenden uns nun zu einer wichtigen Anwendung der Theorie der Invarianten endlicher Gruppen, dem \emph{Formenproblem}. Darunter verstehen wir die Aufgabe, die Variablen $x_{1}$, $x_{2}$, $\ldots$, $x_{n}$ als Functionen der Invarianten einer Gruppe $S$ linearer Substitutionen auszudrücken.

Endlich können wir auch noch nach \emph{gebrochenen Invarianten} fragen. Ist
\[
R(x) = \frac{F_{1}(x)}{F_{2}(x)}
\]
ein Quotient zweier Invarianten vom gleichen Grade, so ist $R(x)$ eine rationale Function der Variablen, die durch alle Substitutionen der Gruppe $S$ ungeändert bleibt, also eine absolute Invariante.

Bedeutet $A$ eine Substitution der endlichen Gruppe $S$, so können wir aus $F(x)$ eine neue Function:
\[
F'(x) = F(Ax)
\]
bilden. Wenn wir mit $A$ alle Substitutionen von $S$ durchlaufen, so erhalten wir ein System von Formen, die wir als \emph{conjugirte Formen} bezeichnen. Diese Formen sind homogene Functionen von $n$ Variablen, und es können also nicht mehr als $n$ von einander unabhängige unter ihnen vorkommen.

Wir betrachten nun als Rationalitätsbereich $\Omega$ den Körper, der aus den absoluten Invarianten der Gruppe $S$ und allen rationalen Zahlen besteht. Dann ist jede Function der Variablen $x$, die durch die Substitutionen von $S$ ungeändert bleibt, eine rationale Function der Invarianten mit Coefficienten aus $\Omega$.

Wenn eine Function $J$ nach dem Satze 4. als rationale Function von $\Omega$ darstellen, so kann sich diese nicht ändern, wenn eine der Substitutionen von $S$ auf sie angewendet wird. Ist aber $S'$ ein Theiler von $S$ vom Index $j$, aber keine anderen gestattet, so genügt $J$ einer Gleichung $j$\textsuperscript{ten} Grades, die als eine \emph{Resolvente} des Formenproblems bezeichnet wird.

Die Auflösung der allgemeinen Gleichungen 2\textsuperscript{ten}, 3\textsuperscript{ten} und 4\textsuperscript{ten} Grades sind also auf Formenprobleme von nur einer Dimension zurückgeführt, und die Theorie der Ikosaedergleichung zeigt, wie auch die Gleichung 5\textsuperscript{ten} Grades durch ein Formenproblem gelöst werden kann.


\subsection*{\S. 45. Einfluss relativer Invarianten.}
\addcontentsline{toc}{subsection}{\S. 45. Einfluss relativer Invarianten.}

Die im Vorhergehenden entwickelte Theorie der absoluten Invarianten lässt sich wesentlich vereinfachen und verallgemeinern, wenn man auch relative Invarianten zulässt. Wir haben gesehen, dass neben den absoluten Invarianten auch relative Invarianten existiren, und diese können zur Vereinfachung des Formenproblems benutzt werden.

Ist $r$ eine solche relative Invariante, so wird eine gewisse Potenz von $r$, deren Grad $\varrho$ ein Theiler des Grades $\mu$ von $S$ ist, eine absolute Invariante, und $r$ wird also durch eine reine Gleichung in $\Omega$ bestimmt.

Alle Substitutionen von $S$, durch die $r$ ungeändert bleibt, bilden für sich eine Gruppe $T_{1}$, und die Gruppe $T_{1}$ ist, wie wir in \S.~40 gesehen haben, ein Normaltheiler von $S$.

Ferner aber sehen wir, dass jede Function der Variablen $x$, die durch die Substitutionen der Gruppe $T_{1}$ ungeändert bleibt, rational durch $r$ ausgedrückt werden kann, d.~h. in dem durch Adjunction von $r$ aus $\Omega$ abgeleiteten Körper $\Omega_{1}$ enthalten ist.

Nach \S.~40 nämlich giebt es eine Substitution $E$ in $S$, und eine primitive $\varrho$\textsuperscript{te} Einheitswurzel $\varepsilon$, so dass $r$ durch $E$ in $\varepsilon r$ übergeht, und alle Werthe, die $r$ annehmen kann:
\[
r,\ \varepsilon r,\ \varepsilon^{2}r,\ \ldots,\ \varepsilon^{\varrho-1}r,
\]
erhält man durch Wiederholung der Substitution $E$. Bedeutet ferner $t$ eine Function der $x$, die die Substitutionen der Gruppe $T_{1}$ gestattet, und durch $E$ und seine Potenzen in
\[
t,\ t_{1},\ t_{2},\ \ldots,\ t_{\varrho-1}
\]
übergeht, so gestatten alle diese Functionen gleichfalls die Substitutionen von $T_{1}$, weil $T_{1}$ ein Normaltheiler von $S$ ist. Die Function
\begin{equation}
(\varepsilon^{-h}, t) = t + \varepsilon^{-h}t_{1} + \varepsilon^{-2h}t_{2} + \cdots + \varepsilon^{-(\varrho-1)h}t_{\varrho-1}
\end{equation}
$(h = 0, 1, 2, \ldots \varrho-1)$, die nach Analogie der Lagrange'schen Resolventen (Bd.~I, \S.~164) gebildet ist, nimmt durch Anwendung der Substitution $E$ den Factor $\varepsilon^{h}$ an, und wenn wir also
\[
(\varepsilon^{-h}, t) = \varepsilon^{h}(\varepsilon^{-h}, t)
\]
setzen, so erhalten wir eine Function, die durch $E$ den Factor $\varepsilon^{h}$ annimmt. Diese Function ist also eine relative Invariante der Gruppe $S$, und aus ihnen können wir durch Linearkombination alle anderen relativen Invarianten gewinnen.

Damit ist gezeigt, dass die Kenntnis einer einzigen relativen Invariante $r$ zur vollständigen Reduktion des Formenproblems auf ein einfacheres zurückreicht, und dass die Auflösung des Formenproblems mit der Auflösung einer reinen Gleichung für $r$ und der Auflösung eines einfacheren Formenproblems für die Gruppe $T_{1}$ gleichbedeutend ist.

\end{document}
